\section{Постановка задачи стабилизации}
\label{sec:foundations}

\subsection{Обратная связь по состоянию и по выходу}
\lecture{L2}{3--9}

Состояние $x\in\mathbb R^n$ содержит переменные, достаточные для описания
дальнейшего движения системы. Вход $u$ выбирает регулятор. Для объекта
\[
  \dot x=f(x,u)
\]
задача стабилизации состоит в построении закона управления, при котором
выбранное равновесие замкнутой системы асимптотически устойчиво.
Нужно обеспечить и малость отклонений при малых начальных возмущениях,
и их стремление к нулю со временем.

\textbf{Статическая обратная связь по состоянию} имеет вид
\[
  u=\gamma(x),\qquad \dot x=f\bigl(x,\gamma(x)\bigr).
\]
Слово «статическая» означает отсутствие собственной памяти регулятора.
Это не означает, что управление постоянно: оно изменяется вместе с $x(t)$.
Перед анализом устойчивости необходимо проверить
$f(0,\gamma(0))=0$. Во всех работах измеряется весь вектор состояния,
поэтому именно такая обратная связь непосредственно реализуема.

\textbf{Динамическая обратная связь} содержит собственное состояние $z$:
\[
  u=\gamma(x,z),\qquad \dot z=g(x,z).
\]
Пример собственной памяти --- накопленная ошибка интегрального регулятора.
Анализировать нужно уже объединённую систему $(x,z)$, а не только объект.
Если целевое равновесие регулятора выбрано нулевым, требование стабилизации
формулируется для $(x,z)=(0,0)$.

При \textbf{обратной связи по выходу} измеряется $y=h(x,u)$, а полный $x$
может быть неизвестен. Статический регулятор использует $u=\gamma(y)$;
динамический --- $u=\gamma(y,z)$ и уравнение для $z$, зависящее от доступных
измерений. Если в правой части этого уравнения непосредственно используется
неизмеряемый $x$, это уже не реализуемая обратная связь только по выходу.
В рассматриваемых работах такой проблемы нет: состояние измеряется целиком.

\addition Если $h$ зависит от $u$, подстановка $u=\gamma(y)$ может создать
алгебраическое уравнение относительно $u$. Нужно отдельно проверить его
разрешимость. В работе 3 выходы заданы как $y=h(x)$, поэтому этой
дополнительной трудности нет.

\subsection{Стабилизация ненулевого равновесия}
\lecture{L2}{10--13}

Равновесие --- состояние, которое можно поддерживать постоянным входом:
\begin{equation}
  f(x_{ss},u_{ss})=0.
  \label{eq:equilibrium-general}
\end{equation}
Задать желаемое $x_{ss}$ недостаточно: необходимо найти совместимое $u_{ss}$.
Если равенство~\eqref{eq:equilibrium-general} не выполняется, точка не является
равновесием и обычное требование её асимптотической устойчивости некорректно.

Вводятся отклонения, обозначенные в лекции индексом $\delta$:
\[
  x_\delta=x-x_{ss},\qquad u_\delta=u-u_{ss}.
\]
Для постоянного задания $\dot x_{ss}=0$, поэтому
\[
  \dot x_\delta=f(x_{ss}+x_\delta,u_{ss}+u_\delta)
  \equiv f_\delta(x_\delta,u_\delta),\qquad f_\delta(0,0)=0.
\]
Если синтезирован закон $u_\delta=\gamma(x_\delta)$, в исходных координатах
нужно записать
\begin{equation}
  u=u_{ss}+\gamma(x-x_{ss}).
  \label{eq:equilibrium-feedback}
\end{equation}
В этой формуле слева стоит полный вход $u$, а не отклонение $u_\delta$.

Для выходного задания сначала выбирают равновесие, удовлетворяющее
$h(x_{ss},u_{ss})=y_{ss}$, затем вводят
\[
  y_\delta=y-h(x_{ss},u_{ss})
  =h(x_{ss}+x_\delta,u_{ss}+u_\delta)-h(x_{ss},u_{ss}).
\]
Нулевой выход сам по себе не определяет единственное состояние:
разные внутренние состояния могут давать одинаковый $y$. Именно поэтому
в работе 3 понадобится анализ нуль-динамики.

\addition В третьей системе работы 4 при $x=0$ первая производная равна
$\dot x_1=\cos0=1$ независимо от $u$. Начало координат стабилизировать нельзя.
Выбрано совместимое равновесие $x_{ss}=(0,1,0,0)^T$, $u_{ss}=0$.
На графике нужно смотреть на $x_2-1$, а не требовать $x_2\to0$.

\subsection{Стабилизация по линейному приближению}
\lecture{L2}{14--17}

Пусть начало координат --- равновесие и $f$ непрерывно дифференцируема.
Матрицы линейного приближения определяются в точке равновесия:
\[
  A=\left.\frac{\partial f}{\partial x}\right|_{x=0,u=0},\qquad
  B=\left.\frac{\partial f}{\partial u}\right|_{x=0,u=0}.
\]
Тогда $\dot x=Ax+Bu$ описывает только первый порядок разложения исходного
объекта. Выбирают $u=-Kx$ так, чтобы $A-BK$ была гурвицевой:
\[
  \operatorname{Re}\lambda_i(A-BK)<0\quad\text{для всех }i.
\]
Производная правой части нелинейной замкнутой системы в нуле равна $A-BK$.
Из теоремы о линеаризации следует локальная экспоненциальная устойчивость
исходной системы. Слово «локальная» существенно: отброшенные нелинейности
могут стать определяющими вдали от равновесия.

\addition Пара $(A,B)$ управляема, если
\[
  \rank\begin{pmatrix}B&AB&\cdots&A^{n-1}B\end{pmatrix}=n.
\]
Тогда линейной обратной связью можно назначить все полюса. Для одной лишь
стабилизации достаточно стабилизируемости: все неуправляемые моды должны
уже иметь отрицательные действительные части. Неуправляемую неустойчивую
моду обратная связь по этому входу устранить не сможет.

Пример различия локального и глобального вывода:
$\dot x=x^3+u$, $u=-x$ дают $\dot x=-x+x^3$.
Линеаризация $\dot x=-x$ устойчива, но при $x>1$ правая часть положительна.
Регулятор стабилизирует нуль локально; глобального результата нет.

\subsection{Маятник и заданный угол}
\lecture{L2}{18--25}

В лекции рассматривается модель
\[
  \ddot\theta=-a\sin\theta-b\dot\theta+cT,
  \qquad a>0,\quad b\ge0,\quad c>0.
\]
Для постоянного угла $\theta=\delta$ скорость и ускорение равны нулю.
Отсюда находятся удерживающий момент и отклонения:
\[
  T_{ss}=\frac{a}{c}\sin\delta,\qquad
  x_1=\theta-\delta,\quad x_2=\dot\theta,\quad u=T-T_{ss}.
\]
В этих координатах
\[
  \dot x_1=x_2,\qquad
  \dot x_2=-a\bigl[\sin(x_1+\delta)-\sin\delta\bigr]-bx_2+cu.
\]
Обратите внимание на скобки: $\sin(x_1+\delta)-\sin\delta$ обращается в нуль
при $x_1=0$. Замена этого выражения на $\sin x_1+\delta-\sin\delta$
изменила бы модель.

Матрицы и матрица замкнутой системы равны
\[
  A=\begin{pmatrix}0&1\\-a\cos\delta&-b\end{pmatrix},\qquad
  B=\begin{pmatrix}0\\c\end{pmatrix},\qquad K=\begin{pmatrix}k_1&k_2\end{pmatrix},
\]
\[
  A-BK=\begin{pmatrix}0&1\\-(a\cos\delta+ck_1)&-(b+ck_2)\end{pmatrix}.
\]
Характеристический полином имеет вид
\[
  p(s)=s^2+(b+ck_2)s+a\cos\delta+ck_1.
\]
Для второго порядка гурвицевость эквивалентна положительности обоих
коэффициентов после старшего:
\[
  k_1>-\frac{a\cos\delta}{c},\qquad k_2>-\frac bc.
\]
Итоговый закон записывается через измеряемые угол и скорость:
\begin{equation}
  T=\frac{a}{c}\sin\delta-k_1(\theta-\delta)-k_2\dot\theta.
\end{equation}
Первое слагаемое удерживает равновесие, второе возвращает угол к заданию,
третье создаёт демпфирование. Даже для верхнего положения, где
$\cos\delta<0$, достаточно сильная позиционная обратная связь делает
линейное приближение устойчивым. Гарантия пока остаётся локальной.

\addition Если задан желаемый полином $s^2+d_1s+d_0$ с $d_0,d_1>0$, то
$k_1=(d_0-a\cos\delta)/c$, $k_2=(d_1-b)/c$.
Так выводится выбор коэффициентов, а не просто угадывается знак $K$.

\subsection{Линейный динамический регулятор по выходу}
\lecture{L2}{26--28}

Для $y=h(x)$ линейный выход имеет вид $y=Cx$, где
$C=\left.\partial h/\partial x\right|_0$.
Предполагаются стабилизируемость $(A,B)$ и обнаруживаемость $(A,C)$.
Обнаруживаемость означает, что все ненаблюдаемые моды уже устойчивы.
Линейный динамический регулятор лекции записан так:
\[
  \dot z=Fz+Gy,\qquad u=Lz+My.
\]
После подстановки в объект анализируется расширенная система
\begin{equation}
  \begin{pmatrix}\dot x\\\dot z\end{pmatrix}
  =\begin{pmatrix}A+BMC&BL\\GC&F\end{pmatrix}
   \begin{pmatrix}x\\z\end{pmatrix}.
\end{equation}
Требуется гурвицевость всей блочной матрицы. Устойчивость одной только
матрицы $F$ ничего не говорит об устойчивости взаимодействия регулятора
с объектом. При гладкости нелинейного объекта гурвицевость расширенного
линейного приближения даёт локальную устойчивость полного замкнутого контура.

\addition Наблюдатель даёт понятную реализацию этого подхода:
\[
  \dot z=Az+Bu+H(y-Cz),\qquad u=-Kz.
\]
Здесь $z$ --- оценка состояния; $H$ выбирается так, чтобы $A-HC$ была
гурвицевой. Ошибка $x-z$ удовлетворяет $\frac{d}{dt}(x-z)=(A-HC)(x-z)$.
Полюса линейного контура состоят из полюсов $A-BK$ и $A-HC$.
В лабораторных 3--5 наблюдатель не нужен, поскольку всё состояние измеримо.

\subsection{Статическая ошибка и интегральное управление}
\lecture{L2}{29--46}

Удерживающее воздействие $T_{ss}=a\sin\delta/c$ требует точных $a,c$.
Если вместо них использованы номинальные $a_0,c_0$, получается закон
\[
  T=\frac{a_0}{c_0}\sin\delta-k_1(\theta-\delta)-k_2\dot\theta.
\]
В установившемся режиме действительного объекта выполняется
\begin{equation}
  a\sin\theta_{ss}
  =c\frac{a_0}{c_0}\sin\delta-ck_1(\theta_{ss}-\delta).
  \label{eq:static-error-pendulum}
\end{equation}
Из этого уравнения в общем случае не следует $\theta_{ss}=\delta$.
Устойчивая система может остановиться с ненулевой статической ошибкой:
устойчивость равновесия и точность задания являются разными требованиями.

На схемах лекции в объекте стоят два интегратора: ускорение интегрируется
в скорость, скорость --- в угол. Позиционная ошибка и скорость возвращаются
с отрицательными знаками через $k_1,k_2$. Канал $a_0\sin\delta/c_0$ подаёт
номинальный удерживающий момент напрямую и сам по себе ошибки модели
не компенсирует.

В иллюстрации лекции коэффициент при $\sin\theta$ в объекте увеличивается
вдвое, задание равно $45^\circ$, а коэффициенты выбраны
$K=\begin{pmatrix}3a/c&a\end{pmatrix}$.
На графике без дополнительного интегратора
угол устанавливается примерно на $36^\circ$. С интегральным каналом
установившийся угол достигает $45^\circ$, хотя переходный процесс имеет
большее перерегулирование. Поэтому улучшение точности не означает
автоматического улучшения всех показателей переходного процесса.

Для устранения постоянной ошибки вводится
\[
  e=\theta-\delta,\qquad z=\int_0^t e(\tau)\,d\tau,\qquad \dot z=e.
\]
Закон лекции содержит интегральное слагаемое с единичным коэффициентом:
\begin{equation}
  T=T_{ss,0}-k_1e-k_2\dot\theta-z,
  \qquad T_{ss,0}=\frac{a_0}{c_0}\sin\delta.
\end{equation}
Если замкнутая система сходится к постоянному равновесию, условие
$\dot z=0$ требует $e=0$. Состояние интегратора при этом становится
$z_{ss}=T_{ss,0}-a\sin\delta/c$: оно накапливает недостающую постоянную
компенсацию. В номинальной модели $z_{ss}=0$.

\addition Интегратор нужно включать в анализ устойчивости. Для отклонений
$(e,\dot\theta,z-z_{ss})$ характеристический полином равен
\[
  p(s)=s^3+(b+ck_2)s^2+(a\cos\delta+ck_1)s+c.
\]
Кроме положительности коэффициентов, критерий Гурвица третьего порядка
требует
\[
  (b+ck_2)(a\cos\delta+ck_1)>c.
\]
Следовательно, механическое добавление интегратора к любому устойчивому
регулятору может сделать контур неустойчивым. Сначала выбирают коэффициенты,
затем проверяют расширенную систему.

Интегральная компенсация постоянной ошибки не равнозначна точному
подавлению произвольного неизвестного $h(t)$ из работы 5. Для переменного
возмущения потребуется отдельный анализ; лекционный непрерывный регулятор
маятника в работе 5 гарантирует малую окрестность нуля.

\subsection{Что нужно различать перед выбором метода}

\begin{center}
\begin{tabularx}{\textwidth}{@{}lYY@{}}
\toprule
Вопрос & Что проверяется & Типичная ошибка\\
\midrule
Равновесие & $f(x_{ss},u_{ss})=0$ & Требовать стабилизацию недостижимой точки\\
Устойчивость & Движение при начальных отклонениях & Считать достаточным один удачный график\\
Точность & Значение ошибки после переходного процесса & Принимать устойчивость за отсутствие ошибки\\
Измеримость & Какие переменные доступны регулятору & Использовать неизвестный параметр в реальном законе\\
Область гарантии & Локальная область или всё пространство & Называть локальную линеаризацию глобальным методом\\
\bottomrule
\end{tabularx}
\end{center}
